% SPDX-License-Identifier: PMPL-1.0-or-later
% SPDX-FileCopyrightText: 2026 Jonathan D.A. Jewell
%
% The Protocol Factory: Malleable, Ductile Protocol Creation
% via Formally Verified Typed Overlays
%
% Intended venue: arXiv preprint (cs.NI / cs.SE / cs.DC)
% Subsequently: HotNets, NSDI, or ACM Computing Surveys
%
% Fourth paper in the series:
%   1. Groove Protocol (the mechanism)
%   2. Typing the Boundaries (the boundary taxonomy)
%   3. Graduated Trust (the trust spectrum + IPv6T)
%   4. The Protocol Factory (this paper — the bigger vision)

\documentclass[11pt,a4paper]{article}
\usepackage[utf8]{inputenc}
\usepackage[T1]{fontenc}
\usepackage{amsmath,amssymb,amsthm}
\usepackage{listings}
\usepackage{hyperref}
\usepackage{booktabs}
\usepackage{graphicx}
\usepackage{xcolor}
\usepackage[margin=2.5cm]{geometry}

\definecolor{codegreen}{rgb}{0.25,0.49,0.31}
\definecolor{codegray}{rgb}{0.5,0.5,0.5}
\definecolor{codeblue}{rgb}{0.22,0.40,0.72}

\lstset{
  basicstyle=\ttfamily\small,
  keywordstyle=\color{codeblue}\bfseries,
  commentstyle=\color{codegray}\itshape,
  stringstyle=\color{codegreen},
  breaklines=true,
  frame=single,
  captionpos=b
}

\newtheorem{theorem}{Theorem}
\newtheorem{definition}{Definition}
\newtheorem{property}{Property}
\newtheorem{observation}{Observation}

\title{The Protocol Factory:\\
Malleable, Ductile Protocol Creation\\
via Formally Verified Typed Overlays}

\author{Jonathan D.A. Jewell\\
\texttt{j.d.a.jewell@open.ac.uk}\\
The Open University}

\date{March 2026}

\begin{document}

\maketitle

\begin{abstract}
Creating a new communication protocol today requires years of
standardisation, vendor adoption, and infrastructure deployment.
We observe that the convergence of three recent technologies ---
dependent types (Idris2, 2021), zero-overhead FFI (Zig, 2023), and
ubiquitous IPv6 (2024) --- enables a fundamentally different approach:
\emph{protocol creation as compilation}. A user defines capability types
in a formally verified language, compiles them into a panel (a
distributable application component), and every instance of that panel
immediately speaks a new, proven protocol over standard TCP/IPv6. No
RFC process, no vendor adoption, no infrastructure change. The protocol
is malleable (reshapable without breaking) and ductile (stretchable from
localhost to WAN without breaking). We analyse why no incumbent has built
this despite having the resources, identify the unsolved governance
problems, and speculate on the implications of democratised protocol
creation for the structure of networked systems.
\end{abstract}

\section{The Cost of a Protocol}

Creating a communication protocol is among the most expensive acts in
computing. HTTP/1.1 took 6 years from proposal to RFC (1993--1999).
HTTP/2 took 7 years (2007--2014). QUIC took 9 years (2012--2021). IPv6
has been ``deploying'' for 25 years and is not yet universal.

The cost is not primarily technical. The cost is \emph{coordination}:
drafting specifications, achieving consensus among competing vendors,
implementing in multiple languages, deploying across infrastructure
controlled by independent operators, and waiting for adoption to reach
critical mass.

This cost is justified for foundational protocols. HTTP needed to be
universal. TCP needed to be reliable everywhere. DNS needed to be
consistent globally. The coordination cost is the price of universality.

But not every protocol needs to be universal. A voice protocol between
two instances of the same application does not need IETF approval. A
typed data exchange between panels in a desktop environment does not need
W3C consensus. A sensor protocol between IoT devices in a factory does
not need IEEE standardisation.

For these cases --- and they are the majority --- the coordination cost
is pure waste. The protocol is used by a closed community of endpoints.
The only requirement is that the endpoints agree on the types. And if the
endpoints are compiled from the same type definitions, agreement is not a
social process. It is a compile step.

\section{The Protocol Factory}

\begin{definition}[Protocol Factory]
A \emph{protocol factory} is a system in which defining a capability
type, compiling it into a distributable component, and distributing
that component constitutes the complete act of protocol creation.
No external standardisation, infrastructure change, or vendor
coordination is required.
\end{definition}

The factory has four stages, corresponding to the OPSM lifecycle
(Mint, Provision, Configure, Harness):

\begin{enumerate}
\item \textbf{Mint}: define capability types in Idris2. The type checker
  proves the interface is consistent.
\item \textbf{Provision}: compile to Zig FFI, generating C headers. The
  ABI is now callable from any language.
\item \textbf{Configure}: embed in a panel with a Groove manifest
  declaring the capability. The panel is distributable.
\item \textbf{Harness}: distribute the panel. Every instance speaks the
  new protocol. Discovery is automatic (Groove). Wire format is GRV6
  over standard TCP/IPv6.
\end{enumerate}

The entire process --- from type definition to running protocol --- takes
hours, not years. The resulting protocol has stronger correctness
guarantees (dependent type proofs) than most IETF RFCs (which rely on
prose specifications and interoperability testing).

\subsection{Malleable Protocols}

\begin{definition}[Malleability]
A protocol is \emph{malleable} if its types can be reshaped --- fields
added, constraints tightened, interfaces extended --- without breaking
existing connections. Groove's structural subtyping enables this: a
provider offering more than a consumer requires is a valid match.
\end{definition}

Traditional protocols are brittle. Adding a field to HTTP requires a
new RFC. Removing a field breaks clients. Groove protocols are malleable
because matching is structural: if the provider's type is a supertype of
the consumer's requirement, they compose. Adding capabilities to a
manifest never breaks existing consumers.

\subsection{Ductile Protocols}

\begin{definition}[Ductility]
A protocol is \emph{ductile} if it can be stretched from localhost to
WAN without changing its type definitions. The same types, the same
proofs, the same GRV6 frames work on loopback, LAN, and the open
internet.
\end{definition}

GRV6 frames are TCP payload. TCP works identically on localhost and
across the internet. The types in the frame header don't know or care
about the network path. A protocol designed and tested on loopback
stretches to global scale without modification.

\section{Why Incumbents Haven't Built This}

Every company with the engineering resources to build a protocol factory
has a structural incentive not to.

\begin{table}[h]
\centering
\small
\begin{tabular}{lp{3cm}p{5cm}}
\toprule
Company & Resources & Why They Won't \\
\midrule
Google & Chrome (3B users), Protobuf, gRPC, QUIC
  & Removes Google as intermediary. Their revenue depends on
    intermediation (ads, cloud, data). \\
Apple & Safari, Swift, App Store
  & Undermines App Store control. Peer-to-peer protocols bypass
    Apple's 30\% commission. \\
Meta & React Native, Thrift, WhatsApp
  & Undermines platform lock-in. Users communicating via typed
    panels don't need Instagram. \\
Microsoft & Edge, TypeScript, Azure
  & Undermines Azure. Direct peer connections reduce cloud compute
    demand. \\
Amazon & AWS SDK, Ion
  & Undermines AWS entirely. The protocol factory makes cloud
    intermediaries optional. \\
\bottomrule
\end{tabular}
\caption{Incumbent incentive misalignment}
\label{tab:incumbents}
\end{table}

The protocol factory is \emph{disintermediating by design}. Every
company whose revenue depends on being an intermediary has a structural
reason to not build it. This is not conspiracy; it is economics.

\subsection{Why Now}

Three technologies converged in 2021--2024:

\begin{enumerate}
\item \textbf{Idris2} (2021): dependent types with quantitative type
  theory. The first language that provides both formal proofs AND
  practical compilation to efficient code. Previous dependent type
  systems (Coq, Agda) were proof assistants, not systems languages.

\item \textbf{Zig} (2023--2024): zero-overhead FFI with no runtime, no
  GC, comptime evaluation, and C ABI compatibility. The first systems
  language that can match C's performance while being safe enough to
  trust at network boundaries.

\item \textbf{IPv6 ubiquity} (2024): link-local addresses on every
  interface, enabling zero-configuration discovery. Before widespread
  IPv6, local discovery required Bonjour, mDNS infrastructure, or
  manual configuration.
\end{enumerate}

The Venn diagram of ``knows Idris2'' $\cap$ ``knows Zig'' $\cap$
``cares about protocol design'' is vanishingly small. The tools
converged, but the communities haven't overlapped until now.

\section{The Typed Overlay Network}

When two instances of the same panel exchange GRV6 frames, they form a
single edge in a typed overlay network. When a third instance connects,
another edge forms. The overlay grows with every installation.

\begin{property}[Infrastructure-Free Growth]
The typed overlay requires zero infrastructure to grow. No relays, no
trackers, no DHT, no registry. Each panel instance is a node. Each
GRV6 connection is an edge. The internet is the wire.
\end{property}

This is unique among overlay networks:

\begin{table}[h]
\centering
\small
\begin{tabular}{llll}
\toprule
Overlay & Preserves & Infrastructure & Growth Mechanism \\
\midrule
Tor & Anonymity & Relays & Volunteers run relays \\
BitTorrent & Availability & Trackers/DHT & Users share files \\
IPFS & Content addressing & DHT & Nodes pin content \\
Signal & Confidentiality & Servers & Users install app \\
\textbf{Typed Overlay} & \textbf{Type safety} & \textbf{None}
  & \textbf{Users install panels} \\
\bottomrule
\end{tabular}
\caption{Overlay networks compared}
\end{table}

\subsection{Cross-Panel Interoperability}

The overlay is not limited to instances of the same panel. Any two
panels that import the same Idris2 ABI package for a shared capability
are Level 4 (Proven) to each other for that capability.

This means:
\begin{itemize}
\item A chat panel and a file-sharing panel that both import the
  \texttt{TextChannel} ABI can exchange text messages with proven
  type safety.
\item A game server panel and a voice panel that both import the
  \texttt{VoiceTransport} ABI can carry voice with proven type safety.
\item Neither panel needs to know the other exists in advance. Groove
  discovery finds the match. The ABI package is the shared proof root.
\end{itemize}

The typed overlay is open. Publishing an ABI package to a registry
is an invitation for any application to join the overlay for that
capability. Proliferation is adoption.

\section{Broadcast: The Typed Discovery Layer}

Previous work \cite{boundaries2026} identified broadcast protocols as a
boundary where types cannot survive. We revise this assessment.

Broadcast is not connectionless chaos. It is the \emph{discovery layer}
of a typed system. The Groove Protocol already performs typed broadcast:
a system announces its capability manifest via multicast or port probing.
Receivers filter by structural type match. Matches upgrade to connections.

The architecture is two primitives:

\begin{enumerate}
\item \textbf{Typed Broadcast Filter}: embeds type hashes in broadcast
  payloads (mDNS TXT records, BLE advertising data, WiFi vendor IEs).
  Receivers filter by type hash before upgrading to a connection.

\item \textbf{Typed Frame Router}: once a connection is established,
  provides proven frame-level translation with zero-copy splice.
\end{enumerate}

The first primitive handles one-to-many discovery. The second handles
one-to-one communication. Together, they cover the full lifecycle:
broadcast $\rightarrow$ discover $\rightarrow$ connect $\rightarrow$
communicate.

\section{Encryption: Types at the Tunnel Endpoints}

TLS Application-Layer Protocol Negotiation (ALPN, RFC~7301) provides a
standard mechanism for type negotiation during the TLS handshake:

\begin{lstlisting}[caption={GRV6 ALPN negotiation}]
ClientHello:
  ALPN: ["grv6", "h2", "http/1.1"]

ServerHello:
  ALPN: "grv6"

-> Tunnel carries GRV6 typed frames
-> Types negotiated BEFORE encryption starts
-> No inspection of encrypted traffic needed
\end{lstlisting}

This means types exist at both tunnel endpoints without existing inside
the tunnel. The tunnel is a trusted channel between proven endpoints.
The ALPN token is the groove: a standard shape that both sides recognise.

\section{Unsolved Problems}

We are honest about what remains unsolved:

\subsection{Schema Evolution}

When an ABI package is updated, existing panels have the old type hash.
New panels have the new hash. They cannot communicate.

This is the most critical unsolved problem. Possible approaches:
\begin{itemize}
\item \textbf{Backward-compatible evolution}: like Protobuf's field
  numbering. New fields are optional. Old fields are never removed.
  The type hash covers the \emph{core} fields only.
\item \textbf{Version ranges}: GRV6 carries a capability version.
  Receivers accept a range of versions, not a single hash.
\item \textbf{Dual-hash}: the frame carries both a \texttt{type\_hash}
  (exact) and a \texttt{compat\_hash} (backward-compatible subset).
  Matching on either is valid.
\end{itemize}

None of these are implemented. This is future work.

\subsection{Governance of the ABI Registry}

Who curates the canonical ABI packages? If centralised, the registry
is a single point of control (the ``npm problem''). If decentralised,
how do you prevent fragmentation (a thousand incompatible ``voice''
definitions)?

We do not propose a solution. We observe that this is a governance
problem, not a technical one, and that governance is harder than code.

\subsection{Malicious Capabilities}

The type system proves that data is well-formed. It does not prove that
data should have been sent. A panel with a proven ``file-read''
capability can exfiltrate documents with mathematical correctness. Formal
verification proves the mechanism works; it says nothing about whether
the mechanism should be used.

Capability signing, publisher reputation, and user consent are necessary
but are social mechanisms, not type-level guarantees.

\subsection{Traffic Analysis}

GRV6 type hashes are consistent per-type. An observer who knows the ABI
can correlate traffic by hash without decrypting. TLS ALPN mitigates
this for the negotiation phase. Per-session hash salting mitigates it
for the data phase. Neither is currently implemented.

\section{Speculative Implications}

\subsection{The End of Protocol Scarcity}

If creating a protocol costs years and millions of dollars, protocols
are scarce. Scarce protocols are controlled by those who can afford the
cost. This concentrates power in standards bodies and large companies.

If creating a protocol costs hours and is free, protocols are abundant.
Abundant protocols are created by anyone with a problem to solve. This
distributes power to individuals and small teams.

We do not claim this is unambiguously good. Abundance creates
fragmentation. Fragmentation creates confusion. But the alternative ---
scarcity controlled by incumbents --- creates stagnation. The history of
the internet suggests that abundance, despite its chaos, produces more
innovation than scarcity.

\subsection{Protocols as Libraries}

In the protocol factory model, a protocol is not a specification
document. It is a \emph{library}: an importable package of type
definitions that, when compiled into an application, gives that
application the ability to speak a new language.

This inverts the traditional relationship between protocols and
applications. Traditionally, protocols are fixed and applications
conform to them. In the factory model, applications define their
protocols by choosing which ABI packages to import. The protocol
emerges from the application, not the other way around.

\subsection{The Parallel Protocol Layer}

The typed overlay does not replace TCP/IP. It runs inside TCP/IP. It is
a parallel layer:

\begin{lstlisting}[caption={The parallel protocol layer}]
Traditional:    Application -> HTTP -> TLS -> TCP -> IPv6 -> Wire
With GRV6:      Application -> GRV6 -> HTTP -> TLS -> TCP -> IPv6 -> Wire
                               ^^^^
                               Typed payload frame (105 bytes)
                               Type hash, cap hash, provenance
                               Invisible to all layers below
\end{lstlisting}

The GRV6 layer is invisible to HTTP, TLS, TCP, and IPv6. It is
processed only by the application endpoints. The internet infrastructure
carries it without knowing it exists --- and this is by design. The
internet's job is to carry bytes faithfully. GRV6's job is to make those
bytes meaningful.

\section{Conclusion}

The protocol factory is not a product to be launched. It is an
observation about what becomes possible when dependent types, zero-
overhead FFI, and ubiquitous IPv6 converge. The result is that protocol
creation --- historically one of the most expensive acts in computing
--- becomes as cheap as writing and distributing a library.

The implications are speculative. We do not know whether abundant
protocols lead to interoperability or fragmentation. We do not know
whether the typed overlay will grow or remain niche. We do not know
whether incumbents will adopt, compete with, or undermine the factory
model.

What we do know is this: the tools exist, the architecture works, the
proofs hold, the tests pass on localhost, and localhost TCP is WAN TCP.
The typed overlay is not a proposal. It is running code.

Whether the world wants it is a question we cannot answer with a proof.

\bibliographystyle{plain}
\begin{thebibliography}{10}

\bibitem{groove2026}
J.~D.~A.~Jewell.
\newblock Groove: Formally Verified System Composition via Dependent Types and
  Linear Resources.
\newblock arXiv preprint, 2026.

\bibitem{boundaries2026}
J.~D.~A.~Jewell.
\newblock Typing the Boundaries: Preserving Type Safety Across Every System
  Crossing Point.
\newblock arXiv preprint, 2026.

\bibitem{trust2026}
J.~D.~A.~Jewell.
\newblock Graduated Trust: A Spectrum of Type Safety from Localhost to the Open
  Internet.
\newblock arXiv preprint, 2026.

\bibitem{brady2021}
E.~Brady.
\newblock \emph{Idris 2: Quantitative Type Theory in Practice}.
\newblock In ECOOP, 2021.

\bibitem{rfc7301}
S.~Friedl, A.~Popov, A.~Langley, and E.~Stephan.
\newblock Transport Layer Security (TLS) Application-Layer Protocol Negotiation
  Extension.
\newblock RFC 7301, 2014.

\end{thebibliography}

\end{document}
